\documentclass[11pt,a4paper]{article}
\usepackage[utf8]{inputenc}
\usepackage[T1]{fontenc}
\usepackage[spanish,es-minimal,es-noshorthands]{babel}
\usepackage{lmodern}
\usepackage[sfdefault,lining]{FiraSans}
\usepackage[a4paper,top=2.3cm,bottom=2.2cm,left=1.8cm,right=1.8cm,headheight=1.3cm,headsep=0.4cm,footskip=1cm]{geometry}
\usepackage{amsmath,amssymb}
\usepackage{icomma}
\usepackage{graphicx}
\usepackage[table]{xcolor}
\usepackage{tikz}
\usepackage{fontawesome5}
\usepackage{ccicons}
\usepackage{fancyhdr}
\usepackage{lastpage}
\usepackage{tasks}
\usepackage{needspace}
\usepackage[most]{tcolorbox}
\usepackage{array}
\usepackage[hidelinks]{hyperref}
\AtBeginDocument{\hypersetup{pdftitle={Ficha de ejercicios de la unidad 1: Números racionales. 3º ESO},pdfsubject={Matemáticas 3º ESO},pdfkeywords={números enteros, fracciones, números racionales, decimales, fracción generatriz}}}

\definecolor{navy}{HTML}{1B2748}
\definecolor{accent}{HTML}{E8603A}
\definecolor{teal}{HTML}{1F9E90}
\definecolor{soft}{HTML}{EEF1F8}
\definecolor{muted}{HTML}{5D6685}

\setlength{\parindent}{0pt}
\setlength{\parskip}{3pt}
\settasks{label=\alph*), label-format=\bfseries\color{navy}, label-width=1.3em, item-indent=1.8em, column-sep=0.6em, after-item-skip=2pt, before-skip=3pt, after-skip=4pt}

% ---------- cabecera y pie ----------
\newcommand{\llicencia}{\href{https://creativecommons.org/licenses/by-nc-sa/4.0/deed.es}{\ccbyncsa}\ \ Licencia Creative Commons Atribución-NoComercial-CompartirIgual 4.0 Internacional (CC BY-NC-SA 4.0)}
\pagestyle{fancy}
\fancyhf{}
\fancyhead[L]{\raisebox{-0.35\height}{\includegraphics[height=0.95cm]{logo-gva.png}}}
\fancyhead[C]{\small\color{muted}Ficha de ejercicios · Unidad 1: Números racionales · 3\textsuperscript{o} ESO}
\fancyhead[R]{\raisebox{-0.35\height}{\includegraphics[height=0.95cm]{logo-vermellar.png}}}
\fancyfoot[L]{\scriptsize\color{muted}\llicencia}
\fancyfoot[R]{\small\color{muted}\thepage}
\renewcommand{\headrulewidth}{0.4pt}
\renewcommand{\headrule}{\color{muted!40}\hrule height \headrulewidth}
\fancypagestyle{primera}{\fancyhf{}\renewcommand{\headrulewidth}{0pt}
  \fancyfoot[L]{\scriptsize\color{muted}\llicencia}
  \fancyfoot[R]{\small\color{muted}\thepage}}

% ---------- niveles y ejercicios ----------
\newcommand{\estrelles}[1]{{\color{accent}\ifcase#1\or\faStar\faStar[regular]\faStar[regular]\or\faStar\faStar\faStar[regular]\or\faStar\faStar\faStar\fi}}
\newcounter{exer}
\newcommand{\ex}[1]{\par\needspace{4\baselineskip}\medskip\stepcounter{exer}\noindent{\color{navy}\bfseries\theexer.}\ \ {\small\estrelles{#1}}\ \ }

\newcommand{\seccio}[2]{\par\needspace{14\baselineskip}\bigskip
  \begin{tcolorbox}[colback=navy,colframe=navy,coltext=white,arc=2mm,boxrule=0pt,left=3mm,right=3mm,top=1.2mm,bottom=1.2mm]
  {\large\bfseries #1\quad #2}
  \end{tcolorbox}\vspace{-2pt}}

% ---------- recta real ----------
\newcommand{\rectabuida}{\begin{tikzpicture}[x=0.42cm,baseline=0]
  \draw[-latex,muted] (-6.7,0)--(6.8,0);
  \foreach \i in {-6,...,6}{\draw[muted] (\i,0.12)--(\i,-0.12) node[below=-1pt,font=\scriptsize,text=muted]{$\i$};}
  \end{tikzpicture}}
% \rectaint{x1}{x2}{estilo izquierda}{estilo derecha}: estilo = ple, buit, inf
\newcommand{\extrem}[2]{\ifx#2p\fill[accent] (#1,0) circle (0.13);\else\ifx#2b\filldraw[fill=white,draw=accent,line width=1.2pt] (#1,0) circle (0.13);\fi\fi}
\newcommand{\rectaint}[4]{\begin{tikzpicture}[x=0.42cm,baseline=0]
  \draw[-latex,muted] (-6.7,0)--(6.8,0);
  \foreach \i in {-6,...,6}{\draw[muted] (\i,0.12)--(\i,-0.12) node[below=-1pt,font=\scriptsize,text=muted]{$\i$};}
  \draw[accent,line width=3pt] (#1,0)--(#2,0);
  \ifx#3i\fill[accent] (#1-0.25,0)--(#1+0.15,0.2)--(#1+0.15,-0.2)--cycle;\else\extrem{#1}{#3}\fi
  \ifx#4i\fill[accent] (#2+0.25,0)--(#2-0.15,0.2)--(#2-0.15,-0.2)--cycle;\else\extrem{#2}{#4}\fi
  \end{tikzpicture}}

\newcolumntype{C}[1]{>{\centering\arraybackslash}p{#1}}
\newcolumntype{M}[1]{>{\centering\arraybackslash}m{#1}}

\def\fmneed{14}
% =====================================================================
% Capçalera FentMates de les fitxes: un color per a cada curs.
% Els colors són els de la pàgina inicial (LEVELS a assets/js/materials.js);
% el logo ve de _eines/logo/pdf/ (python3 _eines/logo/genera.py).
%
% Ús, al final del preàmbul de la unitat:
%   % =====================================================================
% Capçalera FentMates de les fitxes: un color per a cada curs.
% Els colors són els de la pàgina inicial (LEVELS a assets/js/materials.js);
% el logo ve de _eines/logo/pdf/ (python3 _eines/logo/genera.py).
%
% Ús, al final del preàmbul de la unitat:
%   % =====================================================================
% Capçalera FentMates de les fitxes: un color per a cada curs.
% Els colors són els de la pàgina inicial (LEVELS a assets/js/materials.js);
% el logo ve de _eines/logo/pdf/ (python3 _eines/logo/genera.py).
%
% Ús, al final del preàmbul de la unitat:
%   \input{../comu/fentmates.tex}
%   \fmcurs{1eso}
% i, al començament del document:
%   \capcalera{Matemàtiques}{1r ESO}{Fitxa d'exercicis · Unitat 2}{Divisibilitat}
% Opcional abans de \input: \def\fmneed{14} (espai mínim abans de cada secció).
% =====================================================================
\makeatletter
\def\fm@logodir{../../logo/pdf}
\providecommand{\fmneed}{10}
\def\fm@cursnom{}

% clau del curs, color i color del text damunt del color
\newcommand{\fm@curs}[3]{\@namedef{fm@color@#1}{#2}\@namedef{fm@text@#1}{#3}}
\fm@curs{1eso}{E6B84A}{navy}
\fm@curs{2eso}{3B82F6}{white}
\fm@curs{3eso}{FF7A50}{white}
\fm@curs{4eso-a}{A78BFA}{white}
\fm@curs{4eso-b}{6D28D9}{white}
\fm@curs{1batx-mi}{9B1C1C}{white}
\fm@curs{1batx-ccss}{EF6B6B}{white}
\fm@curs{2batx-mii}{1E6B47}{white}
\fm@curs{2batx-ccss}{2BB3A3}{white}

\newcommand{\fmcurs}[1]{%
  \@ifundefined{fm@color@#1}{\PackageError{fentmates}{Curs desconegut: #1}{}}{}%
  \edef\fm@tmp{\noexpand\definecolor{curs}{HTML}{\@nameuse{fm@color@#1}}%
               \noexpand\colorlet{curstext}{\@nameuse{fm@text@#1}}}\fm@tmp
  \colorlet{cursfosc}{curs!72!black}%
  \def\fm@logo{\fm@logodir/vertical-#1.pdf}%
  \def\fm@marca{\fm@logodir/marca-#1.pdf}}

% ---------- capçalera de la primera pàgina ----------
% #1 matèria (petit), #2 curs, #3 tipus de document i unitat, #4 títol
\newcommand{\capcalera}[4]{%
  \gdef\fm@cursnom{#2}%
  \par\noindent
  \begin{tcolorbox}[colback=curs,colframe=curs,boxrule=0pt,arc=2.2mm,width=3.5cm,height=2.15cm,
      valign=center,left=3mm,right=2mm,top=1mm,bottom=1mm,nobeforeafter,box align=center]
    \color{curstext}{\footnotesize #1}\par\vspace{1pt}{\fontsize{19}{21}\selectfont\bfseries #2\par}%
  \end{tcolorbox}\hspace{5mm}%
  \begin{minipage}[c]{\dimexpr\linewidth-3.5cm-5mm-2.9cm\relax}
    {\color{cursfosc}\small\bfseries\MakeUppercase{#3}}\par\vspace{2pt}
    {\color{navy}\fontsize{20}{23}\selectfont\bfseries #4\par}
  \end{minipage}\hfill
  \begin{minipage}[c]{2.6cm}\raggedleft\includegraphics[width=2.6cm]{\fm@logo}\end{minipage}\par
  \vspace{2pt}}

% ---------- barres de secció del color del curs ----------
\renewcommand{\seccio}[2]{\par\needspace{\fmneed\baselineskip}\bigskip
  \begin{tcolorbox}[colback=curs,colframe=curs,coltext=curstext,arc=2mm,boxrule=0pt,left=3mm,right=3mm,top=1.2mm,bottom=1.2mm]
  {\large\bfseries #1\quad #2}
  \end{tcolorbox}\vspace{-2pt}}

% ---------- estrelles de dificultat del color del curs (a Batxillerat no se'n posen) ----------
\@ifundefined{estrelles}{}{%
  \renewcommand{\estrelles}[1]{{\color{curs}\ifcase#1\or\faStar\faStar[regular]\faStar[regular]\or\faStar\faStar\faStar[regular]\or\faStar\faStar\faStar\fi}}}

% ---------- capçalera i peu de les pàgines ----------
% primera pàgina: logos institucionals al peu
\fancypagestyle{primera}{\fancyhf{}\renewcommand{\headrulewidth}{0pt}
  \fancyfoot[L]{\raisebox{-0.32\height}{\includegraphics[height=0.78cm]{logo-gva.png}}\hspace{5mm}%
    \raisebox{-0.32\height}{\includegraphics[height=0.78cm]{logo-vermellar.png}}}
  \fancyfoot[R]{\scriptsize\color{muted}\href{https://creativecommons.org/licenses/by-nc-sa/4.0/deed.ca}{\ccbyncsa}\ \ CC BY-NC-SA 4.0 · fentmates.org\qquad\small\thepage}}
% la resta: el símbol a l'esquerra i el curs a la dreta, amb el títol al centre com abans
\fancyhead[L]{\raisebox{-0.3\height}{\includegraphics[height=0.62cm]{\fm@marca}}}
\fancyhead[R]{\small\bfseries\color{cursfosc}\fm@cursnom}
\renewcommand{\headrule}{\color{curs!45}\hrule height \headrulewidth}
\makeatother

%   \fmcurs{1eso}
% i, al començament del document:
%   \capcalera{Matemàtiques}{1r ESO}{Fitxa d'exercicis · Unitat 2}{Divisibilitat}
% Opcional abans de \input: \def\fmneed{14} (espai mínim abans de cada secció).
% =====================================================================
\makeatletter
\def\fm@logodir{../../logo/pdf}
\providecommand{\fmneed}{10}
\def\fm@cursnom{}

% clau del curs, color i color del text damunt del color
\newcommand{\fm@curs}[3]{\@namedef{fm@color@#1}{#2}\@namedef{fm@text@#1}{#3}}
\fm@curs{1eso}{E6B84A}{navy}
\fm@curs{2eso}{3B82F6}{white}
\fm@curs{3eso}{FF7A50}{white}
\fm@curs{4eso-a}{A78BFA}{white}
\fm@curs{4eso-b}{6D28D9}{white}
\fm@curs{1batx-mi}{9B1C1C}{white}
\fm@curs{1batx-ccss}{EF6B6B}{white}
\fm@curs{2batx-mii}{1E6B47}{white}
\fm@curs{2batx-ccss}{2BB3A3}{white}

\newcommand{\fmcurs}[1]{%
  \@ifundefined{fm@color@#1}{\PackageError{fentmates}{Curs desconegut: #1}{}}{}%
  \edef\fm@tmp{\noexpand\definecolor{curs}{HTML}{\@nameuse{fm@color@#1}}%
               \noexpand\colorlet{curstext}{\@nameuse{fm@text@#1}}}\fm@tmp
  \colorlet{cursfosc}{curs!72!black}%
  \def\fm@logo{\fm@logodir/vertical-#1.pdf}%
  \def\fm@marca{\fm@logodir/marca-#1.pdf}}

% ---------- capçalera de la primera pàgina ----------
% #1 matèria (petit), #2 curs, #3 tipus de document i unitat, #4 títol
\newcommand{\capcalera}[4]{%
  \gdef\fm@cursnom{#2}%
  \par\noindent
  \begin{tcolorbox}[colback=curs,colframe=curs,boxrule=0pt,arc=2.2mm,width=3.5cm,height=2.15cm,
      valign=center,left=3mm,right=2mm,top=1mm,bottom=1mm,nobeforeafter,box align=center]
    \color{curstext}{\footnotesize #1}\par\vspace{1pt}{\fontsize{19}{21}\selectfont\bfseries #2\par}%
  \end{tcolorbox}\hspace{5mm}%
  \begin{minipage}[c]{\dimexpr\linewidth-3.5cm-5mm-2.9cm\relax}
    {\color{cursfosc}\small\bfseries\MakeUppercase{#3}}\par\vspace{2pt}
    {\color{navy}\fontsize{20}{23}\selectfont\bfseries #4\par}
  \end{minipage}\hfill
  \begin{minipage}[c]{2.6cm}\raggedleft\includegraphics[width=2.6cm]{\fm@logo}\end{minipage}\par
  \vspace{2pt}}

% ---------- barres de secció del color del curs ----------
\renewcommand{\seccio}[2]{\par\needspace{\fmneed\baselineskip}\bigskip
  \begin{tcolorbox}[colback=curs,colframe=curs,coltext=curstext,arc=2mm,boxrule=0pt,left=3mm,right=3mm,top=1.2mm,bottom=1.2mm]
  {\large\bfseries #1\quad #2}
  \end{tcolorbox}\vspace{-2pt}}

% ---------- estrelles de dificultat del color del curs (a Batxillerat no se'n posen) ----------
\@ifundefined{estrelles}{}{%
  \renewcommand{\estrelles}[1]{{\color{curs}\ifcase#1\or\faStar\faStar[regular]\faStar[regular]\or\faStar\faStar\faStar[regular]\or\faStar\faStar\faStar\fi}}}

% ---------- capçalera i peu de les pàgines ----------
% primera pàgina: logos institucionals al peu
\fancypagestyle{primera}{\fancyhf{}\renewcommand{\headrulewidth}{0pt}
  \fancyfoot[L]{\raisebox{-0.32\height}{\includegraphics[height=0.78cm]{logo-gva.png}}\hspace{5mm}%
    \raisebox{-0.32\height}{\includegraphics[height=0.78cm]{logo-vermellar.png}}}
  \fancyfoot[R]{\scriptsize\color{muted}\href{https://creativecommons.org/licenses/by-nc-sa/4.0/deed.ca}{\ccbyncsa}\ \ CC BY-NC-SA 4.0 · fentmates.org\qquad\small\thepage}}
% la resta: el símbol a l'esquerra i el curs a la dreta, amb el títol al centre com abans
\fancyhead[L]{\raisebox{-0.3\height}{\includegraphics[height=0.62cm]{\fm@marca}}}
\fancyhead[R]{\small\bfseries\color{cursfosc}\fm@cursnom}
\renewcommand{\headrule}{\color{curs!45}\hrule height \headrulewidth}
\makeatother

%   \fmcurs{1eso}
% i, al començament del document:
%   \capcalera{Matemàtiques}{1r ESO}{Fitxa d'exercicis · Unitat 2}{Divisibilitat}
% Opcional abans de \input: \def\fmneed{14} (espai mínim abans de cada secció).
% =====================================================================
\makeatletter
\def\fm@logodir{../../logo/pdf}
\providecommand{\fmneed}{10}
\def\fm@cursnom{}

% clau del curs, color i color del text damunt del color
\newcommand{\fm@curs}[3]{\@namedef{fm@color@#1}{#2}\@namedef{fm@text@#1}{#3}}
\fm@curs{1eso}{E6B84A}{navy}
\fm@curs{2eso}{3B82F6}{white}
\fm@curs{3eso}{FF7A50}{white}
\fm@curs{4eso-a}{A78BFA}{white}
\fm@curs{4eso-b}{6D28D9}{white}
\fm@curs{1batx-mi}{9B1C1C}{white}
\fm@curs{1batx-ccss}{EF6B6B}{white}
\fm@curs{2batx-mii}{1E6B47}{white}
\fm@curs{2batx-ccss}{2BB3A3}{white}

\newcommand{\fmcurs}[1]{%
  \@ifundefined{fm@color@#1}{\PackageError{fentmates}{Curs desconegut: #1}{}}{}%
  \edef\fm@tmp{\noexpand\definecolor{curs}{HTML}{\@nameuse{fm@color@#1}}%
               \noexpand\colorlet{curstext}{\@nameuse{fm@text@#1}}}\fm@tmp
  \colorlet{cursfosc}{curs!72!black}%
  \def\fm@logo{\fm@logodir/vertical-#1.pdf}%
  \def\fm@marca{\fm@logodir/marca-#1.pdf}}

% ---------- capçalera de la primera pàgina ----------
% #1 matèria (petit), #2 curs, #3 tipus de document i unitat, #4 títol
\newcommand{\capcalera}[4]{%
  \gdef\fm@cursnom{#2}%
  \par\noindent
  \begin{tcolorbox}[colback=curs,colframe=curs,boxrule=0pt,arc=2.2mm,width=3.5cm,height=2.15cm,
      valign=center,left=3mm,right=2mm,top=1mm,bottom=1mm,nobeforeafter,box align=center]
    \color{curstext}{\footnotesize #1}\par\vspace{1pt}{\fontsize{19}{21}\selectfont\bfseries #2\par}%
  \end{tcolorbox}\hspace{5mm}%
  \begin{minipage}[c]{\dimexpr\linewidth-3.5cm-5mm-2.9cm\relax}
    {\color{cursfosc}\small\bfseries\MakeUppercase{#3}}\par\vspace{2pt}
    {\color{navy}\fontsize{20}{23}\selectfont\bfseries #4\par}
  \end{minipage}\hfill
  \begin{minipage}[c]{2.6cm}\raggedleft\includegraphics[width=2.6cm]{\fm@logo}\end{minipage}\par
  \vspace{2pt}}

% ---------- barres de secció del color del curs ----------
\renewcommand{\seccio}[2]{\par\needspace{\fmneed\baselineskip}\bigskip
  \begin{tcolorbox}[colback=curs,colframe=curs,coltext=curstext,arc=2mm,boxrule=0pt,left=3mm,right=3mm,top=1.2mm,bottom=1.2mm]
  {\large\bfseries #1\quad #2}
  \end{tcolorbox}\vspace{-2pt}}

% ---------- estrelles de dificultat del color del curs (a Batxillerat no se'n posen) ----------
\@ifundefined{estrelles}{}{%
  \renewcommand{\estrelles}[1]{{\color{curs}\ifcase#1\or\faStar\faStar[regular]\faStar[regular]\or\faStar\faStar\faStar[regular]\or\faStar\faStar\faStar\fi}}}

% ---------- capçalera i peu de les pàgines ----------
% primera pàgina: logos institucionals al peu
\fancypagestyle{primera}{\fancyhf{}\renewcommand{\headrulewidth}{0pt}
  \fancyfoot[L]{\raisebox{-0.32\height}{\includegraphics[height=0.78cm]{logo-gva.png}}\hspace{5mm}%
    \raisebox{-0.32\height}{\includegraphics[height=0.78cm]{logo-vermellar.png}}}
  \fancyfoot[R]{\scriptsize\color{muted}\href{https://creativecommons.org/licenses/by-nc-sa/4.0/deed.ca}{\ccbyncsa}\ \ CC BY-NC-SA 4.0 · fentmates.org\qquad\small\thepage}}
% la resta: el símbol a l'esquerra i el curs a la dreta, amb el títol al centre com abans
\fancyhead[L]{\raisebox{-0.3\height}{\includegraphics[height=0.62cm]{\fm@marca}}}
\fancyhead[R]{\small\bfseries\color{cursfosc}\fm@cursnom}
\renewcommand{\headrule}{\color{curs!45}\hrule height \headrulewidth}
\makeatother

\fmcurs{3eso}

\begin{document}
\thispagestyle{primera}

% ---------- encabezado ----------
\capcalera{Matemáticas}{3\textsuperscript{o} ESO}{Ficha de ejercicios · Unidad 1}{Números racionales}
\vspace{10pt}
Nombre y apellidos: \rule{6.6cm}{0.4pt}\hfill Grupo: \rule{1.6cm}{0.4pt}\hfill Fecha: \rule{2.2cm}{0.4pt}
\vspace{8pt}

\begin{tcolorbox}[colback=soft,colframe=soft,arc=2mm,left=3mm,right=3mm,top=1.5mm,bottom=1.5mm]
\small Ejercicios graduados por niveles de dificultad:\quad
\estrelles{1}\ básico \qquad \estrelles{2}\ medio \qquad \estrelles{3}\ avanzado.\\[2pt]
Resuélvelos en el cuaderno mostrando todos los pasos. Simplifica siempre los resultados. Las rectas y las tablas se pueden completar directamente en la ficha.
\end{tcolorbox}

% =====================================================
\seccio{1}{Conjuntos numéricos: naturales, enteros y racionales}

\ex{1} Marca con una cruz los conjuntos a los que pertenece cada número:\par\smallskip
{\renewcommand{\arraystretch}{2.1}\centering
\begin{tabular}{|M{2.6cm}|M{1.3cm}|M{1.3cm}|M{1.3cm}|}
\hline
\rowcolor{soft}\textbf{Número} & $\mathbb{N}$ & $\mathbb{Z}$ & $\mathbb{Q}$ \\ \hline
$15$ & & & \\ \hline
$-8$ & & & \\ \hline
$\dfrac34$ & & & \\ \hline
$0$ & & & \\ \hline
$-\dfrac{12}{3}$ & & & \\ \hline
$0,6$ & & & \\ \hline
$1,\overline{3}$ & & & \\ \hline
$\dfrac{20}{4}$ & & & \\ \hline
$-2,5$ & & & \\ \hline
\end{tabular}\par}

\ex{2} Indica el conjunto numérico más pequeño al que pertenece cada número:
\begin{tasks}(4)
  \task $-\dfrac{36}{9}$
  \task $\dfrac{42}{6}$
  \task $-\dfrac72$
  \task $0,125$
  \task $-(-5)$
  \task $\dfrac{18}{-3}$
  \task $2,\overline{7}$
  \task $(-2)^{2}$
\end{tasks}

\ex{2} Indica si las siguientes afirmaciones son verdaderas o falsas. Justifícalo.
\begin{tasks}(2)
  \task Todo número entero es racional.
  \task Todo número racional es entero.
  \task El número $-3$ es natural.
  \task La resta de dos naturales siempre es un natural.
  \task Entre $\frac13$ y $\frac12$ hay algún número racional.
  \task La división de dos enteros (con divisor no nulo) siempre es un racional.
\end{tasks}

\ex{3} Responde:
\begin{tasks}(1)
  \task Escribe tres números racionales comprendidos entre $\frac14$ y $\frac13$.
  \task ¿Cuántos números enteros hay entre $-3,5$ y $2,1$? Escríbelos.
  \task Escribe un número que sea entero pero no natural, y otro que sea racional pero no entero.
\end{tasks}

% =====================================================
\seccio{2}{Números enteros: valor absoluto, opuesto y orden}

\ex{1} Calcula:
\begin{tasks}(4)
  \task $|-12|$
  \task $|9|$
  \task $\text{op}(-7)$
  \task $\text{op}(15)$
  \task $|\text{op}(-4)|$
  \task $-|-6|$
  \task $\text{op}\big(|-3|\big)$
  \task $|0|$
\end{tasks}

\ex{1} Responde:
\begin{tasks}(1)
  \task Ordena de menor a mayor: \ $-3,\ \ 5,\ \ -8,\ \ 0,\ \ 2,\ \ -1,\ \ -12$.
  \task Ordena de mayor a menor: \ $7,\ \ -7,\ \ 4,\ \ -10,\ \ -2,\ \ 1$.
  \task Representa en la recta los números $-5$, $3$ y $-2$ y también sus opuestos.\par\smallskip
  \rectabuida
\end{tasks}

\ex{2} Completa con el signo $<$ o $>$:
\begin{tasks}(3)
  \task $-5\ \square\ -2$
  \task $|-5|\ \square\ |-2|$
  \task $-10\ \square\ 0$
  \task $\text{op}(3)\ \square\ \text{op}(-3)$
  \task $-(-4)\ \square\ -4$
  \task $|-7|\ \square\ 6$
\end{tasks}

\ex{2} Escribe todos los números enteros $x$ que cumplen:
\begin{tasks}(2)
  \task $|x|=8$
  \task $|x|<4$
  \task $|x|\le 2$ \ y \ $x$ negativo
  \task $-4<x\le 1$
\end{tasks}

\ex{3} Las temperaturas mínimas de una semana en Morella han sido: lunes $-3\,^\circ$C, martes $2\,^\circ$C, miércoles $-5\,^\circ$C, jueves $0\,^\circ$C y viernes $-1\,^\circ$C.
\begin{tasks}(1)
  \task Ordena los días del más frío al más caluroso.
  \task ¿Qué diferencia hay entre la temperatura más alta y la más baja?
  \task ¿Qué día la temperatura tuvo el mayor valor absoluto?
\end{tasks}

% =====================================================
\seccio{3}{Suma y resta de números enteros}

\ex{1} Calcula:
\begin{tasks}(3)
  \task $(+7)+(-4)$
  \task $(-6)+(-9)$
  \task $(-12)+(+5)$
  \task $(+3)-(-8)$
  \task $(-4)-(+6)$
  \task $(-2)-(-2)$
\end{tasks}

\ex{1} Quita los paréntesis y calcula:
\begin{tasks}(2)
  \task $8-(-3)+(-5)$
  \task $-4+(-7)-(-10)$
  \task $-(-6)-(+9)+2$
  \task $15+(-20)-(-5)-(+3)$
\end{tasks}

\ex{2} Calcula agrupando los positivos y los negativos:
\begin{tasks}(2)
  \task $-3+8-12+5-1$
  \task $10-15+7-9+11$
  \task $-20+13-6+18-5$
  \task $4-7-8+2-3+10$
\end{tasks}

\ex{2} Calcula, empezando por los paréntesis de dentro:
\begin{tasks}(2)
  \task $7-(3-9)+(-4+1)$
  \task $-(5-8)-(2+6)+4$
  \task $12-\left[5-(3-7)\right]$
  \task $-6-\left[-(2-9)+4\right]-1$
  \task $\left[8-(-3+5)\right]-\left[-4-(6-10)\right]$
\end{tasks}

\ex{3} Resuelve estos problemas:
\begin{tasks}(1)
  \task Un submarino está a $120$ m de profundidad. Sube $45$ m, después baja $30$ m y finalmente sube $70$ m. ¿A qué profundidad se encuentra ahora?
  \task Arquímedes nació en el año $287$ a.\,C. y murió en el año $212$ a.\,C. ¿Cuántos años vivió?
  \task A las 6 h el termómetro marca $-4\,^\circ$C. A las 14 h la temperatura ha subido $11$ grados y a las 22 h ha bajado $9$ grados respecto a las 14 h. ¿Qué temperatura hace a las 22 h?
\end{tasks}

% =====================================================
\seccio{4}{Multiplicación y división de números enteros}

\ex{1} Calcula:
\begin{tasks}(4)
  \task $(-5)\cdot(+6)$
  \task $(-4)\cdot(-9)$
  \task $(+7)\cdot(-8)$
  \task $(-1)\cdot(-1)$
  \task $(-48):(-6)$
  \task $(+54):(-9)$
  \task $(-63):(+7)$
  \task $0:(-5)$
\end{tasks}

\ex{1} Calcula aplicando la regla de los signos:
\begin{tasks}(2)
  \task $(-2)\cdot3\cdot(-5)$
  \task $(-1)\cdot(-2)\cdot(-3)\cdot(-4)$
  \task $(-2)\cdot(-2)\cdot(-2)\cdot5$
  \task $(-3)\cdot(-10)\cdot0\cdot(-7)$
\end{tasks}

\ex{2} Completa el número que falta:
\begin{tasks}(4)
  \task $(-6)\cdot\square=42$
  \task $\square:(-4)=9$
  \task $56:\square=-8$
  \task $(-3)\cdot\square\cdot(-2)=-30$
\end{tasks}

\ex{2} Calcula de izquierda a derecha:
\begin{tasks}(2)
  \task $(-72):(-6):(-3)$
  \task $(-60):(-5)\cdot(-2)$
  \task $100:(-4):(-5)$
  \task $(-2)\cdot(-18):(-4)$
\end{tasks}

\ex{3} Responde:
\begin{tasks}(1)
  \task El producto de dos números enteros es $-36$ y su suma es $0$. ¿Qué números son?
  \task Escribe tres parejas de números enteros cuyo producto sea $-12$.
  \task Sin calcular ninguna potencia, ¿cuál es el resultado de $(-1)^{2025}\cdot(-1)^{2026}$?
\end{tasks}

% =====================================================
\seccio{5}{Operaciones combinadas con números enteros}

\ex{1} Calcula respetando la jerarquía de las operaciones:
\begin{tasks}(2)
  \task $5+3\cdot(-2)$
  \task $-8-12:(-4)$
  \task $4\cdot(-3)-2\cdot(-5)$
  \task $20:(-5)+(-3)\cdot(-6)$
\end{tasks}

\ex{2} Calcula:
\begin{tasks}(2)
  \task $6-2\cdot(5-9)$
  \task $(7-12):(-5)+3$
  \task $-3\cdot(4-6)^{2}$
  \task $18:(-3)\cdot2-(-4)$
  \task $(-2)^{3}+3^{2}-(-1)^{5}$
  \task $-2^{4}+(-2)^{4}$
\end{tasks}

\ex{2} Calcula:
\begin{tasks}(2)
  \task $15-\left[4-2\cdot(3-8)\right]$
  \task $-3\cdot\left[(-2)^{2}-7\right]+10:(-2)$
  \task $\left[(-6)\cdot3+20\right]:(-2)-(-5)$
  \task $2\cdot\left\{5-\left[3-(4-9)\right]\right\}$
\end{tasks}

\ex{3} Calcula:
\begin{tasks}(1)
  \task $(-3)^{2}-4\cdot\left[6-(-2)^{3}:(-4)\right]+(-10):5$
  \task $\left\{-2\cdot\left[3-(-1)^{4}\cdot5\right]+(-3)^{2}\right\}:(-13)$
  \task $25:(-5)^{2}-\left[(-4)\cdot(-3)-2^{4}\right]\cdot(-1)^{3}$
  \task Coloca un paréntesis en $\ 3\cdot4-6:2+1\ $ para que el resultado sea $-2$. ¿Y para que sea $4$?
\end{tasks}

% =====================================================
\seccio{6}{Fracciones y números racionales}

\ex{1} Simplifica hasta obtener la fracción irreducible:
\begin{tasks}(6)
  \task $\dfrac{12}{18}$
  \task $\dfrac{45}{60}$
  \task $\dfrac{28}{35}$
  \task $-\dfrac{40}{56}$
  \task $\dfrac{66}{121}$
  \task $\dfrac{84}{36}$
\end{tasks}

\ex{1} Comprueba si son equivalentes:
\begin{tasks}(4)
  \task $\dfrac35$\ y\ $\dfrac{12}{20}$
  \task $\dfrac46$\ y\ $\dfrac{10}{15}$
  \task $\dfrac78$\ y\ $\dfrac{21}{32}$
  \task $-\dfrac29$\ y\ $\dfrac{6}{-27}$
\end{tasks}

\ex{2} Calcula el valor de $x$ para que las fracciones sean equivalentes:
\begin{tasks}(4)
  \task $\dfrac23=\dfrac{x}{15}$
  \task $\dfrac5x=\dfrac{20}{36}$
  \task $\dfrac{x}{7}=\dfrac{18}{42}$
  \task $\dfrac{12}{16}=\dfrac{9}{x}$
\end{tasks}

\ex{2} Reduce a común denominador y ordena de menor a mayor:
\begin{tasks}(1)
  \task $\dfrac34,\ \ \dfrac56,\ \ \dfrac23,\ \ \dfrac{7}{12}$
  \task $-\dfrac12,\ \ \dfrac35,\ \ -\dfrac34,\ \ 0,\ \ \dfrac{7}{10}$
\end{tasks}

\ex{3} Responde:
\begin{tasks}(1)
  \task Escribe una fracción de denominador $20$ comprendida entre $\frac25$ y $\frac35$.
  \task ¿Qué fracción de una hora son $25$ minutos? Simplifícala.
  \task Encuentra una fracción equivalente a $\frac37$ en la que la suma del numerador y el denominador sea $60$.
\end{tasks}

% =====================================================
\seccio{7}{Operaciones con fracciones}

\ex{1} Calcula y simplifica:
\begin{tasks}(3)
  \task $\dfrac37+\dfrac27$
  \task $\dfrac59-\dfrac89$
  \task $\dfrac12+\dfrac13$
  \task $\dfrac34-\dfrac56$
  \task $2-\dfrac35$
  \task $\dfrac16+\dfrac29-\dfrac13$
\end{tasks}

\ex{1} Calcula y simplifica:
\begin{tasks}(3)
  \task $\dfrac23\cdot\dfrac94$
  \task $\left(-\dfrac56\right)\cdot\dfrac{3}{10}$
  \task $4\cdot\left(-\dfrac78\right)$
  \task $\dfrac35:\dfrac{9}{10}$
  \task $\left(-\dfrac47\right):\left(-\dfrac{2}{21}\right)$
  \task $6:\dfrac34$
\end{tasks}

\ex{2} Calcula y simplifica:
\begin{tasks}(2)
  \task $\dfrac{5}{12}+\dfrac38-\dfrac16$
  \task $-\dfrac{2}{15}-\dfrac{3}{10}+\dfrac16$
  \task $3-\dfrac25-\dfrac{7}{10}$
  \task $\dfrac14-\left(\dfrac23-\dfrac56\right)$
\end{tasks}

\ex{2} Calcula simplificando antes de multiplicar:
\begin{tasks}(2)
  \task $\dfrac{14}{15}\cdot\dfrac{25}{28}$
  \task $\left(-\dfrac{18}{35}\right)\cdot\dfrac{21}{12}\cdot\dfrac{10}{9}$
  \task $\dfrac{12}{25}:\dfrac{18}{35}$
  \task $\left(-\dfrac{8}{27}\right):\left(-\dfrac{16}{9}\right)\cdot\dfrac32$
\end{tasks}

\ex{3} Responde:
\begin{tasks}(1)
  \task ¿Qué fracción hay que sumar a $\frac38$ para obtener $\frac56$?
  \task ¿Por qué fracción hay que multiplicar $\frac49$ para obtener $\frac23$?
  \task La suma de dos números es $\frac{7}{12}$ y uno de ellos es $-\frac14$. ¿Cuál es el otro?
  \task ¿Cuál es el inverso del opuesto de $\frac27$?
\end{tasks}

% =====================================================
\seccio{8}{Potencias de fracciones}

\ex{1} Calcula:
\begin{tasks}(4)
  \task $\left(\dfrac34\right)^{2}$
  \task $\left(-\dfrac25\right)^{3}$
  \task $\left(-\dfrac13\right)^{4}$
  \task $-\left(\dfrac12\right)^{4}$
  \task $\left(\dfrac57\right)^{0}$
  \task $\left(\dfrac23\right)^{-2}$
  \task $\left(-\dfrac14\right)^{-1}$
  \task $\left(\dfrac32\right)^{-3}$
\end{tasks}

\ex{2} Expresa como una sola potencia y, si se puede, calcula:
\begin{tasks}(3)
  \task $\left(\dfrac23\right)^{4}\cdot\left(\dfrac23\right)^{3}$
  \task $\left(\dfrac54\right)^{6}:\left(\dfrac54\right)^{2}$
  \task $\left[\left(-\dfrac12\right)^{3}\right]^{2}$
  \task $\left(\dfrac35\right)^{4}\cdot\left(\dfrac53\right)^{4}$
  \task $\left(\dfrac72\right)^{5}:\left(\dfrac72\right)^{5}$
  \task $\left(\dfrac49\right)^{3}:\left(\dfrac23\right)^{3}$
\end{tasks}

\ex{2} Calcula:
\begin{tasks}(2)
  \task $\left(\dfrac12\right)^{2}+\left(\dfrac12\right)^{3}$
  \task $\left(\dfrac23\right)^{2}-\left(\dfrac13\right)^{2}$
  \task $\left(-\dfrac32\right)^{2}\cdot\left(\dfrac23\right)^{3}$
  \task $\left(-\dfrac12\right)^{3}:\left(\dfrac14\right)^{2}$
\end{tasks}

\ex{3} Simplifica y calcula:
\begin{tasks}(2)
  \task $\dfrac{\left(\frac23\right)^{5}\cdot\left(\frac32\right)^{2}}{\left(\frac23\right)^{2}}$
  \task $\left[\left(\dfrac12\right)^{3}\cdot\left(\dfrac12\right)^{-5}\right]^{2}$
  \task $\left[\left(\dfrac35\right)^{2}\right]^{3}:\left(\dfrac35\right)^{5}$
  \task $\left(-\dfrac23\right)^{4}\cdot\left(-\dfrac32\right)^{3}$
\end{tasks}

% =====================================================
\seccio{9}{Operaciones combinadas con fracciones}

\ex{1} Calcula y simplifica:
\begin{tasks}(2)
  \task $\dfrac12+\dfrac13\cdot\dfrac34$
  \task $\dfrac23-\dfrac16:2$
  \task $\dfrac56\cdot\left(\dfrac12+\dfrac13\right)$
  \task $\left(\dfrac34-\dfrac12\right):\dfrac58$
\end{tasks}

\ex{2} Calcula y simplifica:
\begin{tasks}(2)
  \task $\dfrac23+\dfrac12\cdot\left(\dfrac53-1\right)$
  \task $\left(1-\dfrac34\right)\cdot\left(2+\dfrac23\right)$
  \task $\dfrac35-\dfrac25:\left(\dfrac12-\dfrac13\right)$
  \task $\left(\dfrac12\right)^{2}+\dfrac34:\left(\dfrac32-1\right)$
\end{tasks}

\ex{2} Calcula y simplifica:
\begin{tasks}(2)
  \task $\left[\dfrac23-\left(\dfrac12\right)^{2}\right]:\dfrac56$
  \task $1-\left[\dfrac34-\left(\dfrac12-\dfrac13\right)\right]$
  \task $\left(-\dfrac23\right)^{2}:\dfrac43-\dfrac16$
  \task $\dfrac32-\dfrac23\cdot\left[\dfrac14-\left(1-\dfrac12\right)^{2}\cdot3\right]$
\end{tasks}

\ex{3} Calcula y simplifica:
\begin{tasks}(2)
  \task $\left\{\dfrac12-\left[\dfrac23-\left(\dfrac34-\dfrac16\right)\right]\right\}:\left(\dfrac13\right)^{2}$
  \task $\left(1-\dfrac12\right)^{-2}\cdot\left(\dfrac34-1\right)+\left(-\dfrac12\right)^{3}\cdot8$
  \task $\dfrac{\dfrac23-\dfrac12}{1+\dfrac13}$
  \task $\dfrac{2-\dfrac{1}{1+\frac12}}{\dfrac12+\dfrac13}$
\end{tasks}

% =====================================================
\seccio{10}{Problemas con números racionales}

\ex{1} De los $840$ alumnos de un instituto, $\frac{5}{12}$ viven en el mismo pueblo donde está el centro. ¿Cuántos alumnos viven allí? ¿Cuántos vienen de fuera?

\ex{1} Un ciclista ha recorrido $\frac35$ de una etapa de $165$ km. ¿Cuántos kilómetros ha hecho? ¿Cuántos le faltan?

\ex{1} Los $\frac38$ de una cantidad de dinero son $45$ €. ¿Cuál es la cantidad total?

\ex{2} En una biblioteca, $\frac25$ de los libros son novelas, $\frac14$ son de ciencia y el resto, $840$ libros, son de otros temas. ¿Cuántos libros hay en total? ¿Cuántas novelas?

\ex{2} Una madre reparte $600$ € entre sus tres hijos: el primero recibe $\frac13$, el segundo $\frac25$ y el tercero, el resto. ¿Cuánto dinero recibe cada uno?

\ex{2} ¿Cuántos vasos de $\frac{3}{16}$ de litro se pueden llenar con $6$ botellas de $\frac34$ de litro?

\ex{2} De un depósito lleno de agua se gasta $\frac13$ el lunes y $\frac34$ de lo que quedaba el martes. ¿Qué fracción del depósito queda? Si la capacidad es de $1\,200$ L, ¿cuántos litros quedan?

\ex{3} Una persona gasta $\frac25$ del sueldo en el alquiler, $\frac13$ de lo que le queda en comida y ahorra los $600$ € restantes. ¿Cuál es su sueldo? ¿Cuánto gasta en cada cosa?

\ex{3} Un grifo llena un depósito en $4$ horas y otro, en $12$ horas.
\begin{tasks}(1)
  \task ¿Cuánto tiempo tardan en llenarlo si se abren los dos a la vez?
  \task ¿Y si, además, está abierto un desagüe que vacía el depósito lleno en $6$ horas?
\end{tasks}

\ex{3} Un agricultor vende el primer día $\frac38$ de las sandías que tiene, el segundo día $\frac25$ de las que le quedan y el tercer día, las $60$ que todavía le quedaban. ¿Cuántas sandías tenía?

\ex{3} Una pelota bota y cada vez llega a los $\frac23$ de la altura anterior. Si la dejamos caer desde $81$ cm, ¿a qué altura llega después del tercer bote?

% =====================================================
\seccio{11}{Números decimales: clasificación}

\ex{1} Clasifica cada decimal (exacto, periódico puro, periódico mixto o no periódico) y, si es periódico, indica su parte entera, el anteperiodo y el periodo:
\begin{tasks}(3)
  \task $2,375$
  \task $0,4444\ldots$
  \task $1,2333\ldots$
  \task $3,141\,592\ldots$
  \task $0,121212\ldots$
  \task $5,1060606\ldots$
\end{tasks}

\ex{1} Haz la división, escribe la expresión decimal y clasifícala:
\begin{tasks}(6)
  \task $\dfrac78$
  \task $\dfrac29$
  \task $\dfrac56$
  \task $\dfrac{13}{11}$
  \task $\dfrac{7}{20}$
  \task $\dfrac{11}{15}$
\end{tasks}

\ex{2} Sin hacer la división, indica qué tipo de expresión decimal tiene cada fracción. Justifícalo.
\begin{tasks}(4)
  \task $\dfrac{3}{40}$
  \task $\dfrac{5}{21}$
  \task $\dfrac{7}{30}$
  \task $\dfrac{9}{12}$
  \task $\dfrac{4}{33}$
  \task $\dfrac{11}{18}$
  \task $\dfrac{14}{35}$
  \task $\dfrac{6}{55}$
\end{tasks}

\ex{3} Responde:
\begin{tasks}(1)
  \task Escribe una fracción de denominador $12$ que tenga una expresión decimal exacta.
  \task Ordena de menor a mayor: \ $0,\overline{34}$;\ \ $0,3$;\ \ $0,3\overline{4}$;\ \ $0,\overline{3}$.
  \task Las fracciones $\frac13$ y $\frac16$ dan decimales periódicos. ¿Puede ser exacta su suma? Compruébalo.
\end{tasks}

% =====================================================
\seccio{12}{Paso de decimal a fracción}

\ex{1} Halla la fracción generatriz y simplifícala:
\begin{tasks}(4)
  \task $0,8$
  \task $2,45$
  \task $0,036$
  \task $3,125$
\end{tasks}

\ex{2} Halla la fracción generatriz y simplifícala:
\begin{tasks}(4)
  \task $0,\overline{7}$
  \task $2,\overline{3}$
  \task $0,\overline{18}$
  \task $1,\overline{36}$
\end{tasks}

\ex{2} Halla la fracción generatriz y simplifícala:
\begin{tasks}(4)
  \task $0,1\overline{6}$
  \task $1,2\overline{7}$
  \task $0,0\overline{3}$
  \task $2,1\overline{36}$
\end{tasks}

\ex{3} Responde:
\begin{tasks}(1)
  \task Demuestra que $0,\overline{9}=1$.
  \task Calcula $0,\overline{3}+0,\overline{6}$ pasando antes a fracción.
  \task ¿Qué número es mayor, $0,\overline{45}$ o $0,4\overline{5}$? Compruébalo con las fracciones generatrices.
  \task Halla la fracción generatriz de $0,\overline{142857}$.
\end{tasks}

% =====================================================
\seccio{13}{Operaciones combinadas con fracciones y decimales}

\ex{2} Pasa los decimales a fracción y calcula:
\begin{tasks}(2)
  \task $0,5+0,\overline{3}$
  \task $1,2\cdot0,\overline{6}$
  \task $0,\overline{6}-0,1\overline{6}$
  \task $2,\overline{3}:0,\overline{7}$
\end{tasks}

\ex{2} Calcula y da el resultado en forma de fracción irreducible:
\begin{tasks}(2)
  \task $0,25+1,\overline{3}\cdot0,75$
  \task $\left(0,\overline{3}+0,5\right)^{2}$
  \task $1,5-0,\overline{2}:0,\overline{4}$
  \task $\left(2,4-0,\overline{6}\right)\cdot0,1\overline{6}$
\end{tasks}

\ex{3} Calcula y da el resultado en forma de fracción irreducible:
\begin{tasks}(2)
  \task $\left(1,\overline{6}-0,5\right):\left(0,\overline{3}\right)^{2}$
  \task $\dfrac{0,\overline{3}+0,25}{1,1\overline{6}}$
  \task $\left(0,\overline{9}-0,1\right)\cdot1,\overline{1}$
  \task $2,0\overline{3}\cdot0,6-0,2$
\end{tasks}

\vfill
\begin{center}\color{muted}\small\faCheckCircle[regular]\ \ ¡Ánimo! Si te cuesta algún ejercicio, repasa los apuntes interactivos de la unidad.\end{center}

\end{document}
